\chapter{Análisis de Amenazas mediante STRIDE}
\label{ch:analisis_stride}

En este capítulo se formaliza la descomposición funcional y el modelado preventivo de amenazas bajo el marco STRIDE sobre el sistema más crítico de la UNAP: la plataforma transaccional de matrícula académica.

\section{Metodología de Modelado STRIDE}

La metodología STRIDE (desarrollada originalmente por Praerit Garg y Loren Kohnfelder en Microsoft) se fundamenta en la descomposición funcional de un sistema para responder de forma sistemática a cuatro interrogantes operativas:
\begin{enumerate}[itemsep=2pt,topsep=1.5pt,parsep=0pt]
    \item \textbf{¿En qué estamos trabajando?} Construcción del Diagrama de Flujo de Datos (DFD).
    \item \textbf{¿Qué puede salir mal?} Mapeo de categorías STRIDE sobre los elementos y límites de confianza.
    \item \textbf{¿Qué vamos a hacer al respecto?} Definición de controles preventivos y correctivos.
    \item \textbf{¿Hicimos un buen trabajo?} Evaluación residual del riesgo y verificación de controles.
\end{enumerate}

El análisis modela las seis categorías de amenazas frente a las propiedades clásicas de seguridad de la información:

\small
\begin{longtable}{>{\bfseries}p{2.8cm} >{\bfseries}p{2.6cm} p{9.0cm}}
\toprule
\rowcolor{fg}
\textcolor{white}{Categoría STRIDE} & \textcolor{white}{Propiedad Violada} & \textcolor{white}{Definición en el Contexto del Software} \\
\midrule
\endfirsthead
\toprule
\rowcolor{fg}
\textcolor{white}{Categoría STRIDE} & \textcolor{white}{Propiedad Violada} & \textcolor{white}{Definición en el Contexto del Software} \\
\midrule
\endhead
\bottomrule
\endfoot
S --- Spoofing & Autenticación & Un actor malicioso finge ser un usuario legítimo (estudiante/docente) o suplanta un servicio del sistema. \\
\midrule
T --- Tampering & Integridad & Modificación malintencionada y no autorizada de datos en tránsito, en memoria o en reposo. \\
\midrule
R --- Repudiation & No Repudio & Un usuario realiza una acción transaccional y el sistema carece de evidencia inmutable para adjudicarle la autoría. \\
\midrule
I --- Info Disclosure & Confidencialidad & Exposición de credenciales, tokens, datos personales o metadatos de arquitectura a entidades no autorizadas. \\
\midrule
D --- Denial of Service & Disponibilidad & Saturación o degradación intencional de recursos (red, CPU, memoria, hilos web) para impedir el servicio legítimo. \\
\midrule
E --- Elevation of Priv. & Autorización & Un usuario con permisos limitados elude controles para ejecutar funciones de mayor privilegio administrativo. \\
\bottomrule
\end{longtable}
\normalsize

\section{Selección y Alcance del Sistema Objetivo}

\begin{itemize}[itemsep=1.5pt,topsep=1pt,parsep=0pt]
    \item \textbf{Sistema Evaluado:} Portal y API Transaccional de Matrícula (\url{https://matricula.unap.edu.pe}).
    \item \textbf{Capa de Borde:} Proxy inverso web \texttt{nginx/1.16.1} (Linux) encargado de la terminación TLS y enrutamiento perimetral.
    \item \textbf{Capa de Aplicación:} Servicio web compilado en ASP.NET Core (C\#) ejecutándose sobre Kestrel/Linux, identificado mediante las cookies de sesión \path{.AspNetCore.Cookies} (fragmentadas en chunks) y el proveedor de estado temporal \path{.AspNetCore.Mvc.CookieTempDataProvider}.
    \item \textbf{Capa de Persistencia:} Motor de Base de Datos Académica Central (PostgreSQL / Oracle) alojado en el segmento privado del Data Center OTI.
    \item \textbf{Usuarios del Sistema:} Estudiantes de pregrado/posgrado ($\approx$20,000 usuarios recurrentes), docentes y administradores de la Oficina de Tecnologías de la Información.
\end{itemize}

\section{Diagrama de Flujo de Datos (DFD) y Límites de Confianza}

El modelado STRIDE exige delimitar claramente las fronteras de confianza (\textit{Trust Boundaries}), que representan transiciones lógicas o físicas donde los privilegios cambian y el nivel de control disminuye.

\subsection{Elementos del DFD Nivel 1}
\begin{itemize}[itemsep=1.5pt,topsep=1pt,parsep=0pt]
    \item \textbf{Entidades Externas (EE):} \texttt{EE-01:} Estudiante / Docente (Navegador Web / Cliente HTTP).
    \item \textbf{Procesos (P):}
    \begin{itemize}[itemsep=0.8pt,topsep=0.5pt,parsep=0pt]
        \item \texttt{P-01:} Recepción, Terminación TLS y Reverse Proxy (Nginx 1.16.1).
        \item \texttt{P-02:} Gestión de Sesión y Autenticación (Middleware ASP.NET Core Identity).
        \item \texttt{P-03:} Consulta y Despliegue de Horario/Oferta (Endpoint \path{/matricula/alumno/horario}).
        \item \texttt{P-04:} Lógica de Inscripción y Validación de Reglas (Cálculo de créditos, prerrequisitos y vacantes).
    \end{itemize}
    \item \textbf{Almacenes de Datos (D):}
    \begin{itemize}[itemsep=0.8pt,topsep=0.5pt,parsep=0pt]
        \item \texttt{D-01:} Almacén de Estado Temporal / Sesiones en Memoria (\texttt{CookieTempDataProvider}).
        \item \texttt{D-02:} Base de Datos Académica Central (Tablas: Alumnos, Cursos, Actas, Cupos).
    \end{itemize}
    \item \textbf{Flujos de Datos (F):}
    \begin{itemize}[itemsep=0.8pt,topsep=0.5pt,parsep=0pt]
        \item \texttt{F-01:} Petición inicial HTTPS / Envío de credenciales o cookies (Cliente $\rightarrow$ Nginx).
        \item \texttt{F-02:} Petición interna HTTP proxy-pass / Encabezados \texttt{X-Forwarded-*} (Nginx $\rightarrow$ ASP.NET Core).
        \item \texttt{F-03:} Lectura/Escritura de estado de sesión (ASP.NET Core $\leftrightarrow$ D-01).
        \item \texttt{F-04:} Consultas transaccionales SQL / Stored Procedures (ASP.NET Core $\leftrightarrow$ D-02).
        \item \texttt{F-05:} Retorno de payload JSON con horario y estado HTTP 200 (ASP.NET Core $\rightarrow$ Nginx $\rightarrow$ Cliente).
    \end{itemize}
\end{itemize}

\clearpage

\subsection{Esquema Arquitectónico del DFD Nivel 1}

\begin{cmdbox}{Arquitectura DFD Nivel 1 y Límites de Confianza (matricula.unap.edu.pe)}
[ EE-01: Estudiante / Docente ]
               |
               |  F-01: HTTPS / GET /matricula/alumno/horario + Cookies
               v
=== [ LIMITE DE CONFIANZA 1: Internet / Red Publica ] ==================
               |
               v
+----------------------------------------------------------------------+
| P-01: Proxy Inverso (Nginx 1.16.1)                                   |
|   - Terminacion SSL/TLS 1.2                                          |
|   - Enrutamiento interno hacia socket Kestrel                        |
+----------------------------------+-----------------------------------+
                                   |
                                   |  F-02: HTTP Forwarded (Cleartext)
                                   v
+----------------------------------------------------------------------+
| SERVIDOR DE APLICACION ASP.NET CORE                                  |
|                                                                      |
|  +------------------------+      F-03      +----------------------+  |
|  | P-02: Middleware Auth  | <------------> | D-01: Almacen Sesion |  |
|  | (.AspNetCore.Cookies)  |                | (TempDataProvider)   |  |
|  +-----------+------------+                +----------------------+  |
|              | Identidad validada                                    |
|              v                                                       |
|  +------------------------+                +----------------------+  |
|  | P-03: API Controller   | -------------> | P-04: Motor Reglas   |  |
|  | (horario?start=...&end)|                | (Cupos, Cruces, Req) |  |
|  +------------------------+                +----------+-----------+  |
+-------------------------------------------------------|--------------+
                                                        |
=== [ LIMITE DE CONFIANZA 2: DMZ / Red Privada LAN ] ===|===============
                                                        |
                                                        |  F-04: SQL
                                                        v
                                             +-------------------------+
                                             | D-02: BD Central UNAP   |
                                             | (PostgreSQL / Oracle)   |
                                             +-------------------------+
\end{cmdbox}

\clearpage

\section{Evaluación de Riesgos (Criterio DREAD Adaptado)}

Para cuantificar objetivamente la severidad de las amenazas identificadas en el sistema de matrícula, se aplica el modelo DREAD de Microsoft, evaluando cinco dimensiones con puntajes del 1 al 10:
\[
\text{Riesgo Global} = \frac{\text{Damage} + \text{Reproducibility} + \text{Exploitability} + \text{Affected Users} + \text{Discoverability}}{5}
\]

\begin{itemize}[itemsep=3pt,topsep=2pt,parsep=0pt]
    \item \textbf{Damage (Daño potencial):} Perjuicio si el ataque tiene éxito (1 = Mínimo; 10 = Secuestro total de datos o parálisis institucional).
    \item \textbf{Reproducibility (Reproducibilidad):} Facilidad de replicar el ataque una vez descubierto (1 = Muy difícil o circunstancial; 10 = Funciona el 100\% de las veces).
    \item \textbf{Exploitability (Facilidad de explotación):} Destreza y herramientas requeridas (1 = Habilidades avanzadas de kernel/0-day; 10 = Petición HTTP simple con script público).
    \item \textbf{Affected Users (Usuarios afectados):} Proporción de la universidad impactada (1 = Un usuario individual; 10 = Toda la comunidad universitaria).
    \item \textbf{Discoverability (Descubribilidad):} Facilidad para hallar la debilidad (1 = Requiere código fuente; 10 = Visible en cabeceras HTTP públicas).
\end{itemize}

\noindent\textbf{Escala de Severidad Cualitativa:}
\textbf{Bajo:} $1.0 - 4.9$ \quad$\vert$\quad
\textbf{Medio:} $5.0 - 6.9$ \quad$\vert$\quad
\textbf{Alto:} $7.0 - 8.9$ \quad$\vert$\quad
\textbf{Crítico:} $9.0 - 10.0$.

\section{Matriz de Análisis STRIDE del Sistema de Matrícula (UNAP)}

La siguiente matriz documenta las amenazas concretas detectadas a partir de las vulnerabilidades verificadas en el relevamiento perimetral:

\begingroup
\setlength{\tabcolsep}{2.5pt}
\scriptsize
\begin{longtable}{>{\bfseries\centering\arraybackslash}p{0.6cm} p{2.5cm} p{1.4cm} p{2.6cm} >{\centering\arraybackslash}p{1.4cm} p{5.0cm}}
\toprule
\rowcolor{fg}
\textcolor{white}{Cat.} & \textcolor{white}{Amenaza Identificada} & \textcolor{white}{Comp.} & \textcolor{white}{Vulnerabilidad Base} & \textcolor{white}{DREAD / Nivel} & \textcolor{white}{Controles de Mitigación Propuestos} \\
\midrule
\endfirsthead
\toprule
\rowcolor{fg}
\textcolor{white}{Cat.} & \textcolor{white}{Amenaza Identificada} & \textcolor{white}{Comp.} & \textcolor{white}{Vulnerabilidad Base} & \textcolor{white}{DREAD / Nivel} & \textcolor{white}{Controles de Mitigación Propuestos} \\
\midrule
\endhead
\bottomrule
\endfoot
S & Suplantación de Sesión: Robo de cookies \path{.AspNetCore.Cookies} e impersonación de identidad del alumno. & F-01, EE-01, P-02 & Sin binding de sesión a IP/huella TLS; sin revocación forzada tras logout. & \textbf{7.4} \newline \textcolor{orange!90!black}{\textbf{Alto}} \newline \tiny(8,7,6,9,7) & 1. Revocación de tokens JWT / claims en Redis.\newline 2. MFA obligatorio para docentes y condicional alumnos.\newline 3. Binding de cookie a huella criptográfica de conexión. \\
\midrule
T & Manipulación en Tránsito (SSL Stripping): Interceptación en texto plano de parámetros y cookies en Wi-Fi. & F-01, Límite 1 & Ausencia de HSTS (\texttt{Strict-\allowbreak Transport-\allowbreak Security}) en cabeceras web. & \textbf{8.0} \newline \textcolor{orange!90!black}{\textbf{Alto}} \newline \tiny(8,8,7,8,9) & 1. Inyectar directiva HSTS estricta en Nginx (\texttt{max-age=31536000; preload}).\newline 2. Redirección HTTP $\rightarrow$ HTTPS forzada (código 301). \\
\midrule
T & Manipulación de Horario y Cruces: Peticiones alteradas en \texttt{horario?start} para forzar cupos colapsados. & P-03, P-04, F-04 & Validación de reglas en cliente sin doble chequeo atómico en backend. & \textbf{7.4} \newline \textcolor{orange!90!black}{\textbf{Alto}} \newline \tiny(7,9,8,6,7) & 1. Validación transaccional atómica en backend y base de datos (nivel Serializable).\newline 2. Verificación criptográfica de parámetros cliente. \\
\midrule
R & Repudio de Matrícula: Estudiante alega que el sistema no registró su matrícula o notas sin refutación técnica. & P-02, D-02, F-04 & Logs web volátiles; solo registran código 200 sin User-ID ni hash transaccional. & \textbf{5.8} \newline \textcolor{blue!80!black}{\textbf{Medio}} \newline \tiny(6,6,5,7,5) & 1. Bitácora transaccional inmutable (WORM) con User-ID, IP, timestamp NTP y hash del evento.\newline 2. Envío centralizado a servidor Syslog / SIEM externo. \\
\midrule
I & Fuga de Versiones Obsoletas: Atacantes obtienen tecnología exacta para descargar exploits públicos. & P-01, Límite 1 & Cabecera expuesta: \texttt{Server: nginx/1.16.1} (versión 2019 con CVEs acumulados). & \textbf{8.8} \newline \textcolor{orange!90!black}{\textbf{Alto}} \newline \tiny(7,10,8,9,10) & 1. Ocultar firma con \texttt{server\_tokens off;} en Nginx.\newline 2. Actualizar binario de Nginx a rama estable 1.26+.\newline 3. Desplegar reglas de inspección perimetral WAF. \\
\midrule
I & Fuga de Sesión Temporal: Interceptación de cookie \texttt{.AspNetCore.\allowbreak Mvc.\allowbreak CookieTemp\allowbreak DataProvider}. & D-01, F-01 & Cookie temporal no incluye atributo \texttt{Secure} en directiva \texttt{Set-Cookie}. & \textbf{7.0} \newline \textcolor{orange!90!black}{\textbf{Alto}} \newline \tiny(6,7,6,8,8) & 1. Endurecer política de cookies en \texttt{Program.cs}:\newline \texttt{SecurePolicy = CookieSecurePolicy.Always;}\newline \texttt{SameSite = SameSiteMode.Strict;}. \\
\midrule
D & DoS por Agotamiento de Hilos: Colapso de proxy Nginx o cola Kestrel por concurrencia masiva en apertura. & P-01, P-03, D-02 & Sin limitación de tasa (\textit{Rate Limiting}) por IP ni arquitectura elástica. & \textbf{9.0} \newline \textcolor{red!80!black}{\textbf{Crítico}} \newline \tiny(9,9,8,10,9) & 1. Configurar rate limiting en Nginx (\texttt{limit\_req\_zone}).\newline 2. Balanceador de carga con instancias API redundantes.\newline 3. Implementar sala de espera virtual (\textit{Waiting Room}). \\
\midrule
E & Elevación de Privilegios (IDOR): Modificación de IDs de usuario en URL para acceder a rutas de docentes. & P-03, P-04, F-04 & Control de acceso basado en URL sin verificación de claims de rol en backend. & \textbf{7.4} \newline \textcolor{orange!90!black}{\textbf{Alto}} \newline \tiny(9,7,7,8,6) & 1. Autorización basada en políticas de claims (\texttt{[Authorize(Roles = "Docente,Admin")]}).\newline 2. Validación de pertenencia obligatoria entre sesión y recurso. \\
\bottomrule
\end{longtable}
\endgroup

\clearpage

\section{Procedimiento de Implementación de Controles Mitigadores}

Para transformar los hallazgos de la matriz STRIDE en salvaguardas operativas inmediatas dentro de la infraestructura de la UNAP, se establecen tres niveles de acción técnica:

\subsection{A. Hardening Inmediato en Servidor Web Nginx (P-01)}
Se debe modificar el bloque de configuración perimetral en el archivo institucional \texttt{/etc/\allowbreak nginx/\allowbreak sites-available/\allowbreak matricula.\allowbreak unap.\allowbreak edu.\allowbreak pe}:

\begin{cmdbox}{Configuración Nginx de Mitigación STRIDE - matricula.unap.edu.pe}
# Mitigacion STRIDE: Information Disclosure (Ocultamiento de version)
server_tokens off;

# Mitigacion STRIDE: Tampering y Spoofing (HSTS y proteccion de cabeceras)
add_header Strict-Transport-Security "max-age=31536000; includeSubDomains; preload" always;
add_header X-Frame-Options "SAMEORIGIN" always;
add_header X-Content-Type-Options "nosniff" always;
add_header Referrer-Policy "strict-origin-when-cross-origin" always;

# Mitigacion STRIDE: Denial of Service (Proteccion contra rafagas DoS)
limit_req_zone $binary_remote_addr zone=api_matricula:20m rate=15r/s;

server {
    listen 443 ssl http2;
    server_name matricula.unap.edu.pe;

    location / {
        limit_req zone=api_matricula burst=25 nodelay;
        proxy_pass http://127.0.0.1:5000; # Backend Kestrel ASP.NET Core
        proxy_set_header Host $host;
        proxy_set_header X-Real-IP $remote_addr;
        proxy_set_header X-Forwarded-For $proxy_add_x_forwarded_for;
        proxy_set_header X-Forwarded-Proto $scheme;
    }
}
\end{cmdbox}

\subsection{B. Aseguramiento del Manejo de Sesiones en ASP.NET Core (P-02 / D-01)}
En el archivo de inicialización del backend (\texttt{Program.cs}), forzar el endurecimiento de las cookies transaccionales:

\begin{cmdbox}{Endurecimiento de Cookies y Autenticación en Program.cs}
builder.Services.Configure<CookiePolicyOptions>(options =>
{
    // Mitigacion STRIDE: Information Disclosure y Tampering en Cookies
    options.HttpOnly = Microsoft.AspNetCore.CookiePolicy.HttpOnlyPolicy.Always;
    options.Secure = Microsoft.AspNetCore.Http.CookieSecurePolicy.Always;
    options.MinimumSameSitePolicy = Microsoft.AspNetCore.Http.SameSiteMode.Strict;
});

builder.Services.AddAuthentication(CookieAuthenticationDefaults.AuthenticationScheme)
    .AddCookie(options =>
    {
        options.Cookie.Name = "__Host-UNAP-Auth"; // Cookie de prefijo seguro
        options.ExpireTimeSpan = TimeSpan.FromMinutes(20);
        options.SlidingExpiration = true;
    });
\end{cmdbox}

\clearpage

\subsection{C. Control Transaccional en Base de Datos (P-04 / D-02)}
Para mitigar la manipulación de prerrequisitos y vacantes (\textit{Tampering} y \textit{Elevation of Privilege}), toda inscripción debe ejecutarse dentro de un procedimiento almacenado transaccional que valide atómicamente el cupo con bloqueo pesimista:

\begin{cmdbox}{Función SQL Transaccional con Bloqueo FOR UPDATE contra Race Conditions}
-- Mitigacion STRIDE: Tampering (Bloqueo pesimista contra condiciones de carrera)
CREATE OR REPLACE FUNCTION sp_inscribir_asignatura(
    p_alumno_id INT,
    p_curso_id INT,
    p_grupo_id INT
) RETURNS BOOLEAN AS $$
DECLARE
    v_vacantes INT;
    v_inscritos INT;
BEGIN
    SELECT vacantes_totales INTO v_vacantes 
    FROM oferta_academica 
    WHERE grupo_id = p_grupo_id FOR UPDATE;

    SELECT COUNT(*) INTO v_inscritos 
    FROM matricula_detalle 
    WHERE grupo_id = p_grupo_id;

    IF v_inscritos < v_vacantes THEN
        INSERT INTO matricula_detalle(alumno_id, curso_id, grupo_id, fecha_registro)
        VALUES (p_alumno_id, p_curso_id, p_grupo_id, NOW());
        RETURN TRUE;
    ELSE
        RAISE EXCEPTION 'Cupo colapsado: no hay vacantes disponibles';
        RETURN FALSE;
    END IF;
END;
$$ LANGUAGE plpgsql;
\end{cmdbox}
